\section[第1章\q 人生到底该做什么？]{第1章\\ 人生到底该做什么？}
\label{sec:ch1}

周一早上，柔伊同寻常的上班日一样，搭上了纽约地铁L线，手上拿着一杯浓缩咖啡，一边思考着照片。

L线从布鲁克林区向西行，再往南，到达曼哈顿下城区大约需要40分钟。柔伊在最后一站下车。在接下来的40分钟里，柔伊站在诸多乘客中间，依然思考着照片。

那张照片到底想要说什么呢？

到富尔顿转运站时，地铁的门开了。柔伊加入了要转乘去曼哈顿的人潮中。这股人潮把她带到世贸中心下的开放空间，柔伊在灰砖道上停了下来，任由身旁的人从她周围经过。她抬头看着巨大的屋顶，白色钢筋好像巨鹏张开的翅膀，一个从911事件中的灰烬中新生出来的凤凰。

柔伊迈开步伐，继续感受着这一空间的空旷感。意大利大理石有600呎高，使人仿佛置身于一座天主教堂。

那是“纽约之眼”\footnote{The Oculus：世贸中心车站的外观造型。世贸中心车站在911恐攻后被摧毁，后由西班牙建筑师卡拉乔瓦（Santiago Calatrava）设计，展翅的造型有凤凰浴火重生的意境，在2016年被CNN评为九大新地标之一。}。通往这世界上最著名，也最有纪念意义的地方。柔伊每天都要经过两次，一次是去上班的路上，一次是回程，但她却从没有踏进去过。

她走进白色大理石建物的内部，知名的液晶屏幕装饰墙就在左手边，约有一个足球场的长度。这里轮流播放的商业广告和公益广告，她通常不会多看一眼，便转去搭手扶梯。但是今天屏幕上的图像却教柔伊停下了脚步。

屏幕上是一条渔船，还有捕鱼的人及渔网，很像照片上的那条船，那张始终萦绕在她脑中的照片。因为照片上的船不是搁浅在水边坞石上，而是停泊在一处沙漠中央。

还真奇怪，柔伊想。奇怪得让人不安。

就在她瞪着屏幕时，巨大的字幕出现了，是这句话：

{\kaishu 如果你不知道要去哪里，}

{\kaishu 只怕到了哪都不会快乐。}

片刻之后，图像消失了。屏幕开始播放其他广告。

柔伊继续赶去上班。

到达走道尽头，她步上手扶梯，上去两层就是玻璃顶的中庭了。她走出地铁出口，往回折去西街，她工作的杂志社就在前头的大楼里，太阳照着大楼，反光刺着她的眼睛。这栋大楼就是世界贸易中心一号，算是西区最高建物了。每天搭车来这里上班是柔伊的日常，她喜欢在进去办公室前在楼下站一会儿，仰起头来，注视着高耸入云的楼顶。

但是，今天她的思绪却在别处。

{\kaishu 如果你不知道要去哪里，}

{\kaishu 只怕到了哪都不会快乐。}

这应该是一支广告吧，或者是保险公司、汽车销售、旅游APP，具体是哪一个她不记得了。这句广告词是不是洁西卡写的？看似洁西卡的风格，且不管要打广告的重点是什么。但是这个早上，这句广告文案却是要传给柔伊的消息。盯上了她，且不放过。

就像稍早的那张照片，也始终萦绕她的脑海。

柔伊突然想起左手还握着拿铁，赶紧喝一口。咖啡都凉了。

通常她早就过了马路，进入大楼，搭电梯到33楼的办公室了。今天她却没有走寻常的路。穿过西街马路后，她直接右转，来到反思池，那两个方形水池建在双子星大楼的原址，周边黑色的矮石表面铭刻着许多名字。

这是911事件的纪念碑。

她在北边的水池边上停下来，看着池中不断上涌的水。她伸手感受了一下大理石的表面，读了第一排的人名。名字真是多。2001年的9月11日，有数千人罹难于此，那时柔伊正在念研究所。她又望向一条街外的纽约之眼宏伟的支架。

为什么今天一切看起来都不同了？

{\kaishu 如果你不知道要去哪里，}

{\kaishu 只怕到了哪都会不快乐。}

现在柔伊要去哪里？她希望最终到达哪里？

这个问题，她以前真正的思考过吗？

有个男人停下脚步，只是片刻时间，看了看腕表，就又匆匆离开了。柔伊愣了一下，上班要迟到了。

她转身，往世贸中心一号大楼折返回去，但却被什么绊住似的，不能离开。她朝近处的水泥椅走过去，坐了下来，手上拿着冷掉的拿铁。周围是来往的是旅客、当地纽约客，还有车潮。柔伊轻声的问自己：

“我的人生到底该做什么？”
